\documentclass[a4paper, serif, 10pt]{moderncv}

%% moderncv
% color and style
\moderncvcolor{black}
\moderncvstyle{banking}

%% page layout/margins
\usepackage[top=2.5cm, bottom=2.5cm, left=1.5cm, right=1.5cm]{geometry}

\nopagenumbers{}

%% Babel config for Italian language
% ref:
% - https://latex3.github.io/babel/guides/locale-italian.html
\usepackage[italian]{babel}

%% fontspec
\usepackage{fontspec}

\IfFontExistsTF{Linux Libertine}{
  \setmainfont[Ligatures={TeX, Common}, Numbers=OldStyle]{Linux Libertine}
}{}
\IfFontExistsTF{Linux Libertine O}{
  \setmainfont[Ligatures={TeX, Common}, Numbers=OldStyle]{Linux Libertine O}
}{}

%% CTeX
% CJK support, used to show CJK characters in the resume
%
% - heading=false: disable ctex heading
% - punct=kaiming: use kaiming punctuations styles for CJK
% - scheme=plain: use plain scheme, do not override \normalsize font size
% - space=auto: space settings for CJK characters
%
% CJK fonts are resolved per locale language, e.g. Noto Serif/Sans CJK JP
% for Japanese resumes, with fallback to the next region if not installed
%
% ref:
% - http://ctan.mirrorcatalogs.com/language/chinese/ctex/ctex.pdf
\usepackage[UTF8, heading=false, punct=kaiming, scheme=plain, space=auto]{ctex}

\IfFontExistsTF{Noto Serif CJK SC}{
  \setCJKmainfont{Noto Serif CJK SC}
}{}
\IfFontExistsTF{Noto Sans CJK SC}{
  \setCJKsansfont{Noto Sans CJK SC}
}{}

%% line spacing
\usepackage{setspace}
\setstretch{1.125}

%% URL styling - use normal text instead of monospace
\urlstyle{same}

%% Patch \httplink and \httpslink to handle full URLs
% The original moderncv macros always prepend http:// or https:// to URLs.
% This causes issues when the URL already has a protocol (e.g., https://example.com)
% resulting in malformed URLs like https://https://example.com.
% This patch checks if the URL already contains "://" and uses it directly if so.
% Note: We use \str_set:Nx (x-expansion) to expand macros like \@homepage first.
\ExplSyntaxOn
\str_new:N \g_yamlresume_url_str

\RenewDocumentCommand{\httpslink}{O{}m}{
  \str_set:Nx \g_yamlresume_url_str {#2}
  \str_if_in:NnTF \g_yamlresume_url_str {://}
    { \url{#2} }
    { \url{https://#2} }
}

\RenewDocumentCommand{\httplink}{O{}m}{
  \str_set:Nx \g_yamlresume_url_str {#2}
  \str_if_in:NnTF \g_yamlresume_url_str {://}
    { \url{#2} }
    { \url{http://#2} }
}
\ExplSyntaxOff

\name{Giulia Bianchi}{}
\title{Product Manager senior per SaaS B2B e trasformazione digitale}
\phone[mobile]{+39 347 892 4510}
\email{giulia.bianchi@giuliabianchi.it}
\homepage{https://www.giuliabianchi.it}

\address{Via Alessandro Volta 42, Milano, Lombardia, Italia, 20154}{}{}

\extrainfo{{\small \faLinkedin}\ \href{https://www.linkedin.com/in/giuliabianchi-pm}{@giuliabianchi-pm} {} {} {} • {} {} {} 
{\small \faMedium}\ \href{https://medium.com/@giulia.bianchi}{@giulia.bianchi}}

\begin{document}

\maketitle

\section{Informazioni personali}

\cvline{}{Product Manager con oltre otto anni di esperienza nella definizione, nello sviluppo e nella crescita di prodotti digitali SaaS B2B. Lavoro a stretto contatto con team cross-funzionali di prodotto, design, engineering e go-to-market, combinando ascolto del cliente, analisi dei dati e una solida capacità di prioritizzazione.
%
\begin{itemize}
\item Definisco visione, roadmap e OKR allineati agli obiettivi di business
\item Guido discovery, user research, sperimentazione e analisi dei KPI di prodotto
\item Traduco bisogni complessi in requisiti chiari, user story e prototipi validati
\item Coordino stakeholder internazionali in contesti agile, regolamentati e in rapida evoluzione
\end{itemize}}
\section{Formazione}

\cventry{set 2014–lug 2016}
        {Laurea magistrale, Ingegneria Gestionale, Punteggio: 110/110 e lode}
        {Politecnico di Milano}
        {\url{https://www.polimi.it/}}
        {}
        {\begin{itemize}
\item Percorso focalizzato su product management, innovazione e gestione di progetti complessi
\item Tesi sperimentale sull'adozione di piattaforme SaaS nelle PMI italiane
\end{itemize}
\textbf{Corsi}: Strategia e innovazione tecnologica, Marketing industriale, Gestione dei progetti, Economia digitale, Analytics per il business}

\cventry{set 2011–lug 2014}
        {Laurea triennale, Economia e Management, Punteggio: 108/110}
        {Università Bocconi}
        {\url{https://www.unibocconi.it/}}
        {}
        {\textbf{Corsi}: Economia aziendale, Statistica, Sistemi informativi, Ricerche di mercato}
\section{Esperienza professionale}

\cventry{mar 2021–Presente}
        {Senior Product Manager}
        {NexaCloud S.p.A.}
        {\url{https://www.nexacloud.it}}
        {}
        {\begin{itemize}
\item Guido la strategia di prodotto per una piattaforma SaaS B2B di automazione documentale, con un portafoglio da 6 milioni di euro di ricavi annuali.
\item Definisco roadmap trimestrale e OKR in collaborazione con sales, customer success, design e engineering.
\item Ho ridotto del 32\% il time-to-value per i nuovi clienti tramite onboarding guidato e semplificazione dei flussi.
\item Ho introdotto un processo di discovery continuo con oltre 120 interviste a clienti e stakeholder in due anni.
\item Ho coordinato il lancio di tre moduli chiave, contribuendo a un aumento del 18\% nell'adozione della piattaforma.
\end{itemize}
\textbf{Parole chiave}: Roadmap, OKR, SaaS B2B, Discovery, Go-to-market, Data-driven}

\cventry{gen 2018–feb 2021}
        {Product Manager}
        {FinLab Digital S.r.l.}
        {\url{https://www.finlabdigital.it}}
        {}
        {\begin{itemize}
\item Ho gestito il ciclo di vita di un prodotto di pagamenti digitali per PMI, dalla discovery al go-to-market.
\item Ho aumentato del 27\% l'attivazione mensile grazie a un redesign del flusso di onboarding e a campagne in-app.
\item Ho introdotto metriche di prodotto condivise tra business e technology, riducendo i tempi decisionali del 20\%.
\item Ho coordinato un team distribuito di 12 persone tra prodotto, design, engineering e operations.
\end{itemize}
\textbf{Parole chiave}: Pagamenti digitali, Onboarding, Metriche di prodotto, Go-to-market, PMI}

\cventry{set 2016–dic 2017}
        {Product Owner}
        {RetailHub Italia}
        {\url{https://www.retailhubitalia.it}}
        {}
        {\begin{itemize}
\item Ho gestito il backlog e le priorità di una piattaforma e-commerce B2C con oltre 1 milione di visitatori mensili.
\item Ho collaborato con marketing e UX per lanciare un programma di personalizzazione dell'esperienza di acquisto.
\item Ho contribuito a un incremento del 12\% nel tasso di conversione attraverso test A/B e ottimizzazione del funnel.
\end{itemize}
\textbf{Parole chiave}: E-commerce, Backlog, A/B testing, UX, Conversione}
\section{Lingue}

\cvline{Italiano}{Madrelingua o bilingue \hfill \textbf{Parole chiave}: Madrelingua}
\cvline{Inglese}{Competenza professionale completa \hfill \textbf{Parole chiave}: IELTS 8.0, C1 certificato}
\cvline{Francese}{Competenza professionale limitata \hfill \textbf{Parole chiave}: DELF B1}
\section{Competenze}

\cvline{Strategia di prodotto}{Esperto \hfill \textbf{Parole chiave}: Roadmap, OKR, Visione di prodotto, Prioritizzazione, Jobs-to-be-Done}
\cvline{Discovery e ricerca}{Avanzato \hfill \textbf{Parole chiave}: Interviste utente, User testing, Prototipazione, Design thinking}
\cvline{Analytics di prodotto}{Avanzato \hfill \textbf{Parole chiave}: SQL, Mixpanel, Amplitude, Looker, A/B testing}
\cvline{Metodologie Agile}{Esperto \hfill \textbf{Parole chiave}: Scrum, Kanban, SAFe, Backlog management}
\cvline{Leadership e stakeholder}{Avanzato \hfill \textbf{Parole chiave}: Stakeholder management, Mentoring, Comunicazione, Negoziazione}
\cvline{Strumenti}{Avanzato \hfill \textbf{Parole chiave}: Jira, Confluence, Figma, Miro, Notion}
\section{Premi}

\cventry{nov 2022}
        {Osservatorio Innovazione Digitale}
        {Premio Innovazione Digitale}
        {}
        {}
        {Riconoscimento assegnato per il progetto di onboarding intelligente della piattaforma NexaCloud, con un impatto misurabile sulla riduzione del time-to-value e sull'adozione dei clienti.}
\section{Certificazioni}

\cventry{giu 2019}
        {Scrum Alliance}
        {Certificazione Scrum Product Owner (CSPO)}
        {\url{https://www.scrumalliance.org/community/profile/gbianchi}}
        {}
        {}

\cventry{set 2023}
        {Google}
        {Certificazione Google Analytics 4}
        {\url{https://skillshop.withgoogle.com/}}
        {}
        {}
\section{Pubblicazioni}

\cventry{ott 2021}
        {Product discovery nelle PMI italiane: un modello operativo}
        {Economia \& Management}
        {\url{https://www.economiaemanagement.it/articoli/product-discovery-pmi}}
        {}
        {Articolo che presenta un modello di discovery leggero e misurabile per team di prodotto in aziende PMI, basato su casi studio e interviste sul campo.}
\section{Referenze}

\cventry{\emaillink[marco.ferri@nexacloud.it]{marco.ferri@nexacloud.it}}
        {Vicepresidente Prodotto, NexaCloud S.p.A.}
        {Marco Ferri}
        {+39 02 9475 6321}
        {}
        {Giulia ha guidato con grande equilibrio la trasformazione del nostro processo di prodotto, combinando visione strategica, ascolto del cliente e attenzione ai risultati misurabili.}
\section{Progetti}

\cventry{feb 2022–set 2022}
        {Iniziativa di co-creazione con clienti selezionati per validare nuovi moduli della piattaforma.}
        {Programma Beta Aperta}
        {\url{https://www.nexacloud.it/beta-aperta}}
        {}
        {\begin{itemize}
\item Ho reclutato e coordinato 25 clienti beta in tre segmenti di mercato
\item Ho strutturato un ciclo di feedback bisettimanale con product design e engineering
\item Il programma ha generato 14 miglioramenti prioritizzati e due nuove funzionalità candidate
\end{itemize}
\textbf{Parole chiave}: Co-creazione, Beta testing, Customer feedback, Validazione}

\cventry{apr 2023–dic 2023}
        {Progetto di revisione del packaging e del modello di prezzo per il segmento mid-market.}
        {Evoluzione del modello di pricing}
        {}
        {}
        {\begin{itemize}
\item Ho condotto analisi di valore percepito e benchmark competitivo su cinque mercati
\item Ho definito tre nuovi pacchetti con il supporto di sales, finance e customer success
\item Il nuovo modello ha contribuito a un aumento del 9\% nel valore medio del contratto
\end{itemize}
\textbf{Parole chiave}: Pricing, Packaging, Value-based pricing, Mid-market}
\section{Interessi}

\cvline{Tecnologia e innovazione}{Intelligenza artificiale, Prodotti no-code, Open data}
\cvline{Sport e benessere}{Running, Yoga, Trekking}
\cvline{Musica}{Pianoforte, Concerti}
\section{Volontariato}

\cventry{gen 2020–Presente}
        {Mentor di prodotto}
        {Digital Women Italia}
        {\url{https://www.digitalwomen.it/}}
        {}
        {\begin{itemize}
\item Affianco giovani professioniste nella definizione del portfolio e nella preparazione ai colloqui per ruoli di product management
\item Partecipo a incontri mensili e workshop su discovery, priorizzazione e comunicazione con gli stakeholder
\item Ho contribuito a organizzare due edizioni del programma di mentorship
\end{itemize}}
    
\end{document}